\chapter{Higher-order and pattern derivatives}

\begin{orient}
\textbf{What this chapter is for.} A request for $f^{(50)}$, or for $f^{(20)}(0)$, is never an instruction to differentiate fifty times. It is an instruction to find the pattern, and the order was chosen precisely to make the mechanical route impossible and the structural route short. There are exactly four places the pattern can live, and identifying which one is the whole of the work. It can live in a cyclic four-step orbit, as it does for $\sin$ and $\cos$, where differentiation advances a phase by $\pi/2$ and the answer depends only on $n\bmod4$. It can live in a factorial-times-power orbit, as it does for $1/(ax+b)$ and $\ln(ax+b)$ and therefore for every rational function once it has been split into partial fractions. It can live in a Leibniz expansion, where the binomial sum for $(uv)^{(n)}$ has $n+1$ terms of which all but two or three vanish. Or it can live in a Maclaurin coefficient, where $f^{(n)}(0)=n!\,a_{n}$ reads the answer off a series you already know. That last one is the highest-value move in the chapter, and the one a setter builds against when a problem is meant to be hard. Along the way this chapter installs the complex-exponential form that collapses $\eu^{ax}\sin bx$ to a single power, the parity argument that answers half of these problems in one line, Faa di Bruno's formula for composites, and the summation and fractional-order variants. At the end you should treat a large $n$ as a hint rather than a threat.
\end{orient}

\section{Four homes for the pattern}

Nobody differentiates twenty times. The reason is not stamina but growth: each application of the product rule to a two-factor expression doubles the term count, so a fourth derivative of $\eu^{x}\sin x\cos x$ already has more terms than a page holds, and the twentieth has more than there are atoms in the room. Every problem of this kind therefore has a structural solution, and the shape of the problem tells you which one.

To see why, count. Differentiating a sum of $m$ terms, each a product of two factors, produces up to $2m$ terms, and nothing collapses them unless the factors are related. Starting from the two-term product $\eu^{x}\sin x$ the counts run $2,4,8,16,\dots$, so the twentieth derivative has up to $2^{20}$ terms before simplification, which is over a million. Every technique below works by refusing to enter that doubling at all.

\emph{Home one: a cyclic orbit.} Differentiating $\sin(ax+b)$ multiplies by $a$ and advances the phase by $\pi/2$, so after four steps the function returns to itself scaled by $a^{4}$. The orbit is finite, the answer depends on $n\bmod4$, and the entire computation is one reduction modulo four. Any expression that can be linearised into a sum of $\sin(ax+b)$ and $\cos(ax+b)$ lives here, which is why $\sin^{2}x$ and $\cos^{3}x$ must be linearised before anything else happens to them.

\emph{Home two: a factorial-times-power orbit.} Differentiating $(ax+b)^{-1}$ multiplies by $-a$ and raises the power of the denominator by one, so the $n$-th derivative is $(-1)^{n}n!\,a^{n}(ax+b)^{-n-1}$. The orbit is infinite but the pattern is explicit, and $\ln(ax+b)$ joins it after one differentiation. Every rational function lives here once partial fractions have split it, and no rational function lives anywhere else.

\emph{Home three: a Leibniz expansion that truncates.} When $f=uv$ and $u$ is a polynomial of degree $d$, the binomial sum $\sum_{k}\binom nk u^{(k)}v^{(n-k)}$ has $n+1$ terms of which only the first $d+1$ are nonzero, and evaluating at a well chosen point usually kills all but one of those. A twentieth derivative then costs three terms, or one.

\emph{Home four: a Maclaurin coefficient.} At the origin, derivatives and coefficients are the same information in different units, related by $f^{(n)}(0)=n!\,a_{n}$. If you can name the series, you can name the derivative, and naming a series is often possible where finding a derivative pattern is not. This is the most general of the four and the one to reach for when the other three fail to apply, which is exactly the situation a hard problem is built to produce.

\begin{keynote}[The bridge between derivatives and series]
If $f(x)=\sum_{k\geq0}a_{k}x^{k}$ near $0$, then
\[
f^{(n)}(0)=n!\,a_{n}=n!\,[x^{n}]f(x).
\]
Differentiating $n$ times kills every term below $x^{n}$, leaves $n!\,a_{n}$ from the $x^{n}$ term, and sends every term above it to zero at $x=0$. That is the entire content, and it runs in both directions: a hard derivative becomes an easy coefficient, and a hard coefficient becomes an easy derivative. The factor $n!$ is where almost every error in this chapter lives. Two symptoms of the intended route: an order in the teens or higher, and an answer normalised by $n!$ or $\Gamma(n)$, which is the setter cancelling the factorial back out so that the printed answer is small.
\end{keynote}

Two free wins sit above all four homes and should be checked first, because each can end the problem before it starts. The first is parity: at the origin, an odd function has all even derivatives zero and an even function has all odd derivatives zero. The second is the evaluation point: a large part of the expression may vanish at the point in question, and the terms that survive are often a single one.

\begin{center}
\begin{tikzpicture}
\node[dtest] (a) at (0,0) {Is the order $n$ symbolic,\\ or a large specific number?};
\node[dleaf] (L1) at (-3.9,-1.45) {Differentiate directly;\\ Leibniz or Faa di Bruno\\ to organise the terms};
\node[dtest] (b) at (2.9,-1.45) {Is the evaluation\\ point $x=0$?};
\node[dleaf] (L2) at (-0.6,-3.1) {Closed form from three\\ orders, then induction;\\ reduce $n\bmod4$ for trig};
\node[dtest] (c) at (3.6,-3.1) {Is $f$ a polynomial\\ times something else?};
\node[dleaf] (L3) at (1.5,-4.75) {Leibniz: the sum\\ truncates at $\deg$};
\node[dleaf] (L4) at (5.4,-4.75) {$f^{(n)}(0)=n!\,[x^{n}]f$:\\ check parity, then\\ read the coefficient};
\draw[dedge] (a) -- node[dlab,above left]{no} (L1);
\draw[dedge] (a) -- node[dlab,above right]{yes} (b);
\draw[dedge] (b) -- node[dlab,left]{no} (L2);
\draw[dedge] (b) -- node[dlab,right]{yes} (c);
\draw[dedge] (c) -- node[dlab,left]{no} (L3);
\draw[dedge] (c) -- node[dlab,right]{yes} (L4);
\end{tikzpicture}
\end{center}

\section{Closed forms and cyclic patterns}

The elementary families all have closed-form $n$-th derivatives, and there are only four of them worth memorising. Everything else is reduced to these four by an identity, a partial fraction, or a complex exponential.

\tech[n-th derivative closed forms]{$n$-th derivative closed forms and the cyclic trig pattern}{core}
\cue $f^{(n)}$ is wanted symbolically in $n$, or $f^{(n)}(a)$ for a large specific $n$, and $f$ is or can be reduced to one of the elementary families.
\method Establish the pattern from three orders, then prove it by induction. The standard table:
\[
\frac{\mathrm{d}^{n}}{\mathrm{d}x^{n}}\sin(ax+b)=a^{n}\sin\!\left(ax+b+\frac{n\pi}{2}\right),\qquad
\frac{\mathrm{d}^{n}}{\mathrm{d}x^{n}}\eu^{ax}=a^{n}\eu^{ax},
\]
\[
\frac{\mathrm{d}^{n}}{\mathrm{d}x^{n}}\frac{1}{ax+b}=\frac{(-1)^{n}n!\,a^{n}}{(ax+b)^{n+1}},\qquad
\frac{\mathrm{d}^{n}}{\mathrm{d}x^{n}}\ln(ax+b)=\frac{(-1)^{n-1}(n-1)!\,a^{n}}{(ax+b)^{n}} .
\]
The trigonometric entry is periodic with period $4$ in $n$, so reduce $n\bmod4$ before evaluating. Reduce any rational function to a sum of the third form by partial fractions \emph{before} differentiating anything.

\begin{worked}
$f(x)=\sin^{2}x$, and $f^{(7)}(\pi/4)$.

Do not differentiate $\sin^{2}x$. Reduce it first: $\sin^{2}x=\tfrac12-\tfrac12\cos2x$. The constant dies at the first derivative, and the cosine is in the table with $a=2$ and a phase, so for $n\geq1$
\[
f^{(n)}(x)=-\tfrac12\cdot2^{n}\cos\!\left(2x+\frac{n\pi}{2}\right)=-2^{n-1}\cos\!\left(2x+\frac{n\pi}{2}\right).
\]
At $n=7$ and $x=\pi/4$ the argument is $\tfrac\pi2+\tfrac{7\pi}{2}=4\pi$, so $f^{(7)}(\pi/4)=-2^{6}\cos4\pi=-64$. The Cauchy integral formula applied to $\sin^{2}z$ on a circle of radius $0.6$ about $\pi/4$ returns $-64.0000000$.

Now a rational function, where the partial fraction step is not optional. Take $f(x)=\dfrac{1}{x^{2}-5x+6}$ and find $f^{(n)}(0)$. Factor and split:
\[
f(x)=\frac{1}{(x-2)(x-3)}=\frac{1}{x-3}-\frac{1}{x-2},
\]
and each piece is the third table entry with $a=1$. So
\[
f^{(n)}(x)=(-1)^{n}n!\left[\frac{1}{(x-3)^{n+1}}-\frac{1}{(x-2)^{n+1}}\right],
\]
and at $x=0$ the two signs $(-1)^{n}$ and $(-1)^{n+1}$ multiply to $-1$, leaving
\[
f^{(n)}(0)=n!\left[\frac{1}{2^{n+1}}-\frac{1}{3^{n+1}}\right].
\]
At $n=0$ this is $\tfrac12-\tfrac13=\tfrac16=f(0)$, a useful sanity check. At $n=12$ it is $58171.43299$, and contour differentiation returns $58171.43299$. Attempting the same by twelve applications of the quotient rule to $1/(x^{2}-5x+6)$ is not a calculation anybody finishes.

A third specimen, where the pattern starts late. For $f(x)=x\ln x$ the first derivative is $\ln x+1$ and the second is $1/x$, at which point the table takes over: for $n\geq2$,
\[
f^{(n)}(x)=\frac{\mathrm{d}^{\,n-2}}{\mathrm{d}x^{\,n-2}}\frac1x=\frac{(-1)^{n-2}(n-2)!}{x^{n-1}}=\frac{(-1)^{n}(n-2)!}{x^{n-1}} .
\]
At $x=1$ that gives $1,-1,2,-6,24,\dots$ for $n=2,3,4,5,6$, and contour differentiation of $z\ln z$ about $z=1$ returns $1.000$, $-1.000$, $-6.000$ at $n=2,3,5$. The lesson is that a closed form need not be valid from $n=0$; state the range it holds on, and handle the first one or two orders separately.
\end{worked}

\begin{trap}
Getting the logarithm family off by one: the $n$-th derivative of $\ln(ax+b)$ carries $(n-1)!$, not $n!$, because the first differentiation produces $a/(ax+b)$ with no factorial at all and the table for $1/(ax+b)$ then starts from $n-1$. The second standard slip is losing the alternation of signs, which is best avoided by checking $n=1$ against a direct computation every time. The third is differentiating a rational function before splitting it: the quotient rule squares the denominator at every step, and after four rounds the expression is unrecognisable while the partial-fraction form is still two terms.
\end{trap}

\begin{seealsob}Complex exponentials for high derivatives; Maclaurin coefficient extraction; trig identity reduction before high-order differentiation.\end{seealsob}

The induction that justifies the trigonometric entry is worth writing once, because the same three lines justify all four entries and because the phase-shift form is what makes the induction trivial. Suppose $\frac{\mathrm{d}^{n}}{\mathrm{d}x^{n}}\sin(ax+b)=a^{n}\sin\bigl(ax+b+\tfrac{n\pi}{2}\bigr)$. Differentiating once more gives $a^{n+1}\cos\bigl(ax+b+\tfrac{n\pi}{2}\bigr)$, and $\cos\theta=\sin\bigl(\theta+\tfrac\pi2\bigr)$ turns that into $a^{n+1}\sin\bigl(ax+b+\tfrac{(n+1)\pi}{2}\bigr)$, which is the statement at $n+1$. The base case $n=0$ is trivial. The whole force of the phase-shift notation is that differentiation becomes addition of $\pi/2$, so the induction step is an identity rather than a case analysis on $n\bmod4$.

\begin{keynote}[Reduce $n$ modulo $4$, but reduce the function first]
Since $\sin\bigl(\theta+\tfrac{n\pi}{2}\bigr)$ depends only on $n\bmod4$, the derivatives of a sinusoid cycle with period $4$ up to the factor $a^{n}$. But the reduction is only useful after the function has been brought into the form $\sin(ax+b)$ or $\cos(ax+b)$. Powers and products of sinusoids must be linearised first, by $\sin^{2}x=\tfrac12(1-\cos2x)$, $\cos^{3}x=\tfrac14(3\cos x+\cos3x)$, $\sin x\cos x=\tfrac12\sin2x$, and their relatives. A solver who differentiates $\cos^{3}x$ four times without linearising will get the right answer and will take twenty times as long; a solver who tries it eight times will not get the right answer.
\end{keynote}

\begin{keynote}[Induction hygiene]
A conjectured $f^{(n)}$ is worth nothing until it has been checked at both ends. Check the base case at the smallest $n$ for which the formula is claimed, which is often $1$ or $2$ rather than $0$; the $x\ln x$ example above holds only for $n\geq2$, and a solver who asserts it at $n=1$ produces $(-1)!$ and a nonsense answer. Check the step by differentiating the claimed $f^{(n)}$ symbolically and matching it against the claim at $n+1$, not by testing another numerical case. And check the sign at one odd and one even order, since a lost factor of $(-1)^{n}$ is invisible at even orders.
\end{keynote}

One extension of the second home is needed often enough to state separately. Partial fractions do not always produce distinct linear factors; a repeated root leaves a power $(x-c)^{-m}$ with $m\geq2$, and the table entry generalises without difficulty, since each differentiation multiplies by the current exponent and lowers it by one:
\[
\frac{\mathrm{d}^{n}}{\mathrm{d}x^{n}}\,(x-c)^{-m}=(-1)^{n}\,\frac{(m+n-1)!}{(m-1)!}\,(x-c)^{-m-n} .
\]
At $m=1$ this is the table entry, since $(m+n-1)!/(m-1)!=n!$. As a specimen, take $f(x)=\dfrac{2x+1}{(x-1)^{2}(x+2)}$ and find $f^{(n)}(0)$. The partial fractions are $\dfrac{1/3}{x-1}+\dfrac{1}{(x-1)^{2}}-\dfrac{1/3}{x+2}$, checked at $x=0$ where the three pieces give $-\tfrac13+1-\tfrac16=\tfrac12$ against $f(0)=\tfrac12$. Applying the two formulas at $x=0$, with $(0-1)^{-n-1}=(-1)^{n+1}$ and $(0-1)^{-n-2}=(-1)^{n}$,
\[
f^{(n)}(0)=(n+1)!-\frac{n!}{3}-\frac{(-1)^{n}n!}{3\cdot2^{n+1}} .
\]
At $n=0$ that is $1-\tfrac13-\tfrac16=\tfrac12$, and at $n=1,2,3$ it gives $1.75$, $5.25$, $22.125$, all confirmed by contour differentiation. The $(n+1)!$ from the double pole dominates everything else, which is the general lesson: the order of the worst pole sets the growth rate of the derivatives, and the poles nearest the origin set the size of the terms.

Products of an exponential and a sinusoid do not linearise, and their derivatives do not cycle: each differentiation mixes the two terms. The standard escape is to stop treating them as two real functions and start treating them as one complex one, at which point the pattern reappears in its simplest possible form.

\tech[Complex exponentials for high derivatives]{Complex exponentials for $\eu^{ax}\sin(bx)$-type high derivatives}{medium}
\cue A product of a real exponential and a sinusoid, differentiated many times, typically evaluated at $0$.
\method Write the function as the real or imaginary part of a single complex exponential, $\eu^{ax}\sin bx=\operatorname{Im}\eu^{(a+\iu b)x}$ and $\eu^{ax}\cos bx=\operatorname{Re}\eu^{(a+\iu b)x}$. Differentiation of $\eu^{cx}$ is multiplication by $c$, so
\[
\left.\frac{\mathrm{d}^{n}}{\mathrm{d}x^{n}}\bigl(\eu^{ax}\sin bx\bigr)\right|_{x=0}=\operatorname{Im}\bigl[(a+\iu b)^{n}\bigr],
\]
and the power is evaluated in polar form, $(a+\iu b)^{n}=r^{n}\eu^{\iu n\theta}$ with $r=\sqrt{a^{2}+b^{2}}$ and $\theta=\arg(a+\iu b)$. Taking real and imaginary parts commutes with differentiation because $a$ and $b$ are real constants. Keeping $x$ instead of setting it to zero gives the general $n$-th derivative in amplitude-phase form, which is the statement worth memorising:
\[
\frac{\mathrm{d}^{n}}{\mathrm{d}x^{n}}\bigl(\eu^{ax}\sin bx\bigr)=r^{n}\eu^{ax}\sin(bx+n\theta),\qquad
\frac{\mathrm{d}^{n}}{\mathrm{d}x^{n}}\bigl(\eu^{ax}\cos bx\bigr)=r^{n}\eu^{ax}\cos(bx+n\theta).
\]
Each differentiation scales the amplitude by $r$ and advances the phase by $\theta$: the derivatives spiral rather than cycle. When $a=0$ the modulus is $r=b$ and the phase step is $\theta=\pi/2$, the spiral closes after four steps, and the formula degenerates to the cyclic entry of the previous technique. That is the precise sense in which the first home is a special case of this one.

\begin{worked}
$f(x)=\eu^{x}\sin x$, and $f^{(7)}(0)$.

Here $c=1+\iu$, so $r=\sqrt2$ and $\theta=\pi/4$. Then $c^{7}=2^{7/2}\eu^{7\iu\pi/4}$, and $\eu^{7\iu\pi/4}=\tfrac{\sqrt2}{2}(1-\iu)$, so $c^{7}=8-8\iu$ and
\[
f^{(7)}(0)=\operatorname{Im}(8-8\iu)=-8 .
\]
Contour differentiation returns $-8.0000000$. The binomial route, expanding $(1+\iu)^{7}$ into eight terms, gets there too, but the polar form generalises: $f^{(n)}(0)=2^{n/2}\sin(n\pi/4)$ for every $n$, which is the symbolic answer the binomial route never produces.

A second instance with a real part and an unequal pair. For $f(x)=\eu^{2x}\cos3x$ the constant is $c=2+3\iu$ and
\[
f^{(5)}(0)=\operatorname{Re}\bigl[(2+3\iu)^{5}\bigr]=122 ,
\]
which contour differentiation confirms as $122.0000000$. Here the binomial expansion is actually competitive, since $(2+3\iu)^{2}=-5+12\iu$ and squaring again gives $-119-120\iu$, so $(2+3\iu)^{5}=(-119-120\iu)(2+3\iu)=122-597\iu$. Repeated squaring is the practical method for a specific moderate $n$; polar form is the method for symbolic $n$.

A third instance, evaluated away from the origin and with $a<0$, because that is where the branch of $\arg$ starts to matter. For $f(x)=\eu^{-x}\sin x$ the constant is $c=-1+\iu$, so $r=\sqrt2$ and $\theta=\arg(-1+\iu)=3\pi/4$, \emph{not} $\arctan(1/(-1))=-\pi/4$. The amplitude-phase form gives
\[
f^{(n)}(x)=2^{n/2}\eu^{-x}\sin\!\left(x+\frac{3n\pi}{4}\right),
\]
so $f^{(4)}(0.7)=4\eu^{-0.7}\sin(0.7+3\pi)=-4\eu^{-0.7}\sin(0.7)=-1.2796361$, and contour differentiation about $x=0.7$ returns $-1.2796361$. Had the argument been taken as $-\pi/4$, the phase at $n=4$ would have been $-\pi$ rather than $3\pi$, which is the same angle modulo $2\pi$ only by luck; at $n=5$ the two choices differ by $5\pi$ and the sign of the answer flips.
\end{worked}

\begin{trap}
Taking $\operatorname{Re}$ when the function was a sine, or $\operatorname{Im}$ when it was a cosine. Write the identification down before computing anything. The second trap is a wrong branch of $\arg$: for $a<0$ the argument is not $\arctan(b/a)$ but $\arctan(b/a)\pm\pi$, and a sign error there flips the answer for every odd $n$.
\end{trap}

\begin{seealsob}$n$-th derivative closed forms; parity annihilation at the origin; summing all derivatives at a point.\end{seealsob}

Before any of these machines is started, one free check should be run, because it settles a large fraction of problems at the origin without a single differentiation.

\tech{Parity annihilation at the origin}{easy}
\cue A high derivative at $x=0$ of a function visibly assembled from odd and even pieces.
\method If $f$ is odd, every even-order derivative vanishes at $0$; if $f$ is even, every odd-order derivative vanishes at $0$. Determine the parity by the multiplication rules (odd times even is odd, odd times odd is even, and composition with an even inner function is even) before computing anything. The reason is the bridge: an odd function has only odd-degree Maclaurin coefficients, and $f^{(n)}(0)=n!\,a_{n}$.

\begin{worked}
$f(x)=\sinh x\cos x$, and $f^{(10)}(0)$.

$\sinh$ is odd, $\cos$ is even, so $f$ is odd, so all even derivatives vanish at $0$ and $f^{(10)}(0)=0$. Contour differentiation returns $-4\times10^{-9}$, which is zero to the accuracy of the quadrature. The mechanical alternative is a ten-fold product rule with $2^{10}$ terms before cancellation.

The same one-line argument handles orders that no other method could reach: $f(x)=\dfrac{x}{1+x^{2}}$ is odd, so $f^{(2026)}(0)=0$. And it composes with the other routes rather than competing with them: for $g(x)=x^{2}\eu^{-x^{2}}\sin x$, the factors are even, even, odd, so $g$ is odd and every even derivative at $0$ vanishes; the odd ones still need the series, but half the work has already been deleted.
\end{worked}

\begin{trap}
Applying parity at a point other than the centre of symmetry. Oddness of $f$ says nothing whatever about $f^{(10)}(1)$. If the function is symmetric about $x=c$ rather than $0$, substitute $t=x-c$ first and then apply the rule at $t=0$. The other trap is a misremembered parity: $\sinh$ is odd and $\cosh$ is even, matching $\sin$ and $\cos$, and $\tanh$ is odd.
\end{trap}

\begin{seealsob}Maclaurin coefficient extraction; complex exponentials for high derivatives; symmetry kills the term.\end{seealsob}

\begin{keynote}[Parity is the case $k=2$ of a sharper statement]
If $f(x)=g\bigl(x^{k}\bigr)$ for some integer $k\geq2$, then the series of $f$ contains only powers that are multiples of $k$, so
\[
f^{(n)}(0)=0 \text{ unless } k\mid n,\qquad\text{and}\qquad f^{(n)}(0)=n!\,[x^{n/k}]g \text{ when } k\mid n .
\]
Evenness is the case $k=2$ with $g$ arbitrary. The gain grows with $k$: for $f(x)=\sin\bigl(x^{3}\bigr)$, two out of every three derivatives vanish at the origin, and the surviving ones are read straight off the sine series. At $n=15$, $[x^{15}]f=[x^{5}]\sin=1/120$, so $f^{(15)}(0)=15!/120=10897286400$, which contour differentiation confirms. Likewise $\eu^{x^{2}}$ has $f^{(9)}(0)=0$ and $f^{(20)}(0)=20!/10!=670442572800$, since $[x^{20}]\eu^{x^{2}}=[x^{10}]\eu^{x}=1/10!$. Whenever the variable appears only as $x^{k}$, substitute $t=x^{k}$ before doing anything else.
\end{keynote}

\begin{keynote}[The evaluation point is part of the problem]
Parity is the sharpest instance of a more general economy: the point at which the derivative is evaluated usually annihilates most of what you would otherwise compute. In the Leibniz expansion of $x^{2}\eu^{3x}$ below, the general $n$-th derivative has three terms and the value at $0$ has one. In a Faa di Bruno expansion, only the terms whose inner factors survive at the point contribute. The discipline is to carry the evaluation point through the computation rather than saving it for the end: substitute the point as early as the algebra permits, and let it delete terms before you write them down rather than after.
\end{keynote}

\section{Reading the coefficient}

The four homes are not equals. The fourth, the Maclaurin coefficient, is the general one, and the first three are the special cases where the coefficient happens to be available in closed form: a cyclic orbit is a series with period-four coefficients, a factorial-times-power orbit is a geometric series with the factorial divided back out, and a truncating Leibniz sum is a Cauchy product with finitely many terms. That is why a setter who wants a genuinely hard high-derivative problem builds it out of a function whose series you know and whose derivative pattern you do not. Coefficient extraction is accordingly the highest-value technique in the chapter: it is the one that still works when the other three have nothing to say.

Coefficient extraction is only as good as the series you can name, so it is worth having the working set in one place. The following are the coefficients, not the series, because the coefficient is what the technique consumes.

\begin{center}
\footnotesize
\setlength{\tabcolsep}{6pt}
\renewcommand{\arraystretch}{1.32}
\begin{tabular}{@{}lll@{}}
\toprule
\mh{$f(x)$} & \mh{$[x^{n}]f$} & \mh{$f^{(n)}(0)$} \\
\midrule
$\dfrac{1}{1-ax}$ & $a^{n}$ & $n!\,a^{n}$ \\
$\dfrac{1}{(1-x)^{k}}$ & $\dbinom{n+k-1}{k-1}$ & $n!\dbinom{n+k-1}{k-1}$ \\
$\eu^{ax}$ & $\dfrac{a^{n}}{n!}$ & $a^{n}$ \\
$\ln(1+x)$ & $\dfrac{(-1)^{n-1}}{n}$ $(n\geq1)$ & $(-1)^{n-1}(n-1)!$ \\
$\sin x$ & $\dfrac{(-1)^{m}}{(2m+1)!}$ at $n=2m+1$ & $\sin(n\pi/2)$ \\
$\arctan x$ & $\dfrac{(-1)^{m}}{2m+1}$ at $n=2m+1$ & $(-1)^{m}(2m)!$ \\
$(1+x)^{\alpha}$ & $\dbinom{\alpha}{n}$ & $\alpha(\alpha-1)\cdots(\alpha-n+1)$ \\
\bottomrule
\end{tabular}
\end{center}

Two operations extend the table indefinitely. Multiplying by $1/(1-x)$ replaces each coefficient by the partial sum of all coefficients up to it. Multiplying by $x^{m}$ shifts the coefficients right by $m$. Between them, and the Cauchy product, almost every function a setter can build out of the table is reachable.

\tech[Maclaurin coefficient extraction]{Maclaurin coefficient extraction: $f^{(n)}(0)=n!\,[x^{n}]f(x)$}{core}
\cue An absurdly high derivative order at $0$, such as $n=15$, $20$, $100$, $2026$; or an answer normalised by $n!$ or $\Gamma(n)$; or a function built visibly out of series you already know.
\method Expand $f$ as a power series only as far as $x^{n}$, read off the coefficient, and multiply by $n!$. For rational functions get the coefficient by partial fractions plus the geometric family
\[
\frac{1}{1-x}=\sum_{n\geq0}x^{n},\qquad
\frac{1}{(1-x)^{2}}=\sum_{n\geq0}(n+1)x^{n},\qquad
\frac{1}{(1-x)^{k}}=\sum_{n\geq0}\binom{n+k-1}{k-1}x^{n}.
\]
For products, use the Cauchy product: $[x^{n}]\bigl(fg\bigr)=\sum_{j=0}^{n}[x^{j}]f\cdot[x^{n-j}]g$. The rule runs backwards too: a printed series with recognisable coefficients is often an elementary function in disguise, and identifying it answers questions about $f$ at points other than $0$.

\begin{worked}
Let $f(x)=\dfrac{u-ux-x}{(1-x)^{2}}$ with $u$ a constant, and evaluate $\dfrac{f^{(15)}(0)}{\Gamma(15)}$.

The normalisation by $\Gamma(15)=14!$ is the tell. Split the fraction: the numerator is $u(1-x)-x$, so
\[
f(x)=\frac{u}{1-x}-\frac{x}{(1-x)^{2}} .
\]
The first piece contributes $u$ to every coefficient. The second is $x\sum_{n\geq0}(n+1)x^{n}=\sum_{n\geq1}nx^{n}$, contributing $n$ with a minus sign. So $[x^{15}]f=u-15$, and
\[
\frac{f^{(15)}(0)}{\Gamma(15)}=\frac{15!\,(u-15)}{14!}=15u-225 .
\]
At $u=\pi$ that is $15\pi-225=-177.87611$. Contour differentiation of $f$ at $u=\pi$ on a circle of radius $0.5$ gives $f^{(15)}(0)/14!=-177.87611$. The factorials cancelling down to a single factor of $15$ is the whole design of the problem; the answer is small precisely because the setter divided the $15!$ back out.

Now the reverse direction. A function is specified by
\[
f(x)=x+x^{2}-\frac{x^{3}}{6}-\frac{x^{4}}{6}+\frac{x^{5}}{120}+\frac{x^{6}}{120}-\cdots ,
\]
and $f'(6)$ is wanted. The coefficients come in equal pairs $\bigl(1,1\bigr),\bigl(-\tfrac16,-\tfrac16\bigr),\bigl(\tfrac1{120},\tfrac1{120}\bigr)$, which are the coefficients of $\sin x$ repeated. Pairing a series with its own shift is multiplication by $1+x$, so $f(x)=(1+x)\sin x$. Then $f'(x)=\sin x+(1+x)\cos x$ and
\[
f'(6)=\sin6+7\cos6 .
\]
No coefficient was ever extracted; the series was read as a product. This is the direction most often missed.

A rational specimen, to show the geometric family in action. For $f(x)=\dfrac{1}{1-3x+2x^{2}}$, factor the denominator as $(1-x)(1-2x)$ and split:
\[
f(x)=\frac{2}{1-2x}-\frac{1}{1-x}=\sum_{n\geq0}\bigl(2^{n+1}-1\bigr)x^{n},
\]
so $f^{(n)}(0)=n!\bigl(2^{n+1}-1\bigr)$. At $n=10$ that is $3628800\cdot2047=7428153600$, and contour differentiation returns $7.4281536\times10^{9}$. Note that the split was written in the form $1/(1-ax)$ rather than $1/(x-r)$, because that is the form whose geometric expansion is immediate; converting between the two is where signs are lost.

Finally an algebraic specimen, to show that the table's binomial row is not decoration. For $f(x)=(1-4x)^{-1/2}$ the general binomial series gives
\[
[x^{n}]f=\binom{-1/2}{n}(-4)^{n}=\binom{2n}{n},
\]
the standard evaluation of the central binomial coefficient, so
\[
f^{(n)}(0)=n!\binom{2n}{n}=\frac{(2n)!}{n!} .
\]
At $n=8$ that is $40320\cdot12870=518918400$, and contour differentiation on a circle of radius $0.15$ returns $518918400.0$. The radius has to be small here: the nearest singularity is the branch point at $x=1/4$, and a contour outside it computes nothing. That constraint is worth carrying into every numerical check in this chapter.
\end{worked}

\begin{trap}
Reporting the coefficient instead of $n!$ times it. This is the single most punished slip in the whole chapter, and it is invisible in the working: every line is correct until the last one. The countermeasure is to write $f^{(n)}(0)=n!\,a_{n}$ on the page before computing $a_{n}$. The second trap is expanding further than needed: the target is one coefficient, and every term beyond $x^{n}$ is waste that also creates opportunities for error.
\end{trap}

\begin{seealsob}Summing all derivatives at a point; Lagrange inversion and series reversion; the special-function ODE recursion.\end{seealsob}

\begin{keynote}[How large should the answer be]
The coefficients of a function analytic at $0$ satisfy $\abs{a_{n}}\approx C/R^{n}$, where $R$ is the distance from the origin to the nearest singularity of $f$ in the complex plane. Through the bridge that becomes a magnitude estimate for the derivative itself,
\[
\abs{f^{(n)}(0)}\approx \frac{C\,n!}{R^{n}} ,
\]
and it is the cheapest possible check on a long computation. For $f=1/(x^{2}-5x+6)$ the poles are at $2$ and $3$, so $R=2$ and the estimate at $n=12$ is about $12!/2^{13}=58471.9$, against the exact $58171.43$: agreement to half a percent, which is more than enough to catch a dropped factorial or a sign. For $\ln(1+x)/(1-x)$ both singularities sit at distance $1$, so $f^{(20)}(0)$ should be a modest multiple of $20!$, and $1.63\times10^{18}$ against $20!=2.43\times10^{18}$ is exactly that. An answer that misses this estimate by a factor of $n$ or of $n!$ has a located error, not a mysterious one.
\end{keynote}

\subsection{A twentieth derivative, in full}

It is worth doing one of these completely, on a function whose twentieth derivative is genuinely hopeless by any other method, because the contrast is the argument of this chapter. Take
\[
f(x)=\frac{\ln(1+x)}{1-x},\qquad\text{and find } f^{(20)}(0).
\]

The direct route is not merely long, it is impossible in practice. The first derivative is $\dfrac{1}{(1+x)(1-x)}+\dfrac{\ln(1+x)}{(1-x)^{2}}$, already two terms; the second has four, the third has eight, and the logarithm never disappears, so no cancellation ever reduces the count. Nor does the closed-form route work: neither factor is in the elementary table in a way that survives the product, and Leibniz does not truncate because neither factor is a polynomial. Parity is no help either, since $f$ is neither odd nor even.

What is available is that both factors have famous series. From the standard table,
\[
\ln(1+x)=\sum_{k\geq1}\frac{(-1)^{k-1}}{k}x^{k},\qquad \frac{1}{1-x}=\sum_{m\geq0}x^{m},
\]
both valid for $\abs{x}<1$, which is a neighbourhood of $0$ and therefore enough. The Cauchy product of the two says that the coefficient of $x^{n}$ in the product is the sum of $\bigl(\text{coefficient of }x^{k}\text{ in the first}\bigr)\times\bigl(\text{coefficient of }x^{n-k}\text{ in the second}\bigr)$ over all $k$. Every coefficient of the second series is $1$, so the product collapses to a partial sum:
\[
[x^{n}]f=\sum_{k=1}^{n}\frac{(-1)^{k-1}}{k}\cdot1=1-\frac12+\frac13-\cdots+\frac{(-1)^{n-1}}{n} .
\]
Multiplying a series by $1/(1-x)$ is exactly the operation ``replace each coefficient by the sum of all coefficients up to it'', and recognising that saves the Cauchy product entirely.

That partial sum has a closed form. Splitting the alternating sum at $n=2m$ into its positive and negative halves,
\[
\sum_{k=1}^{2m}\frac{(-1)^{k-1}}{k}=\sum_{k=1}^{2m}\frac1k-2\sum_{j=1}^{m}\frac{1}{2j}=H_{2m}-H_{m},
\]
where $H_{m}$ is the $m$-th harmonic number. At $n=20$, so $m=10$,
\[
[x^{20}]f=H_{20}-H_{10}=\frac{155685007}{232792560}=0.66877140 .
\]
The bridge finishes it:
\[
f^{(20)}(0)=20!\,\bigl(H_{20}-H_{10}\bigr)=\frac{2432902008176640000\cdot155685007}{232792560}=1627055289796608000 ,
\]
the division being exact because $232792560=\lcm(1,2,\dots,20)$ divides $20!$, with quotient $10450944000$. Numerically that is $1.6270553\times10^{18}$, and the Cauchy integral formula
\[
f^{(20)}(0)=\frac{20!}{2\pi\iu}\oint_{\abs{z}=0.55}\frac{f(z)}{z^{21}}\dd z
\]
evaluated by the trapezoid rule on $65536$ points returns $1.6270553\times10^{18}$, agreeing to eleven digits.

\begin{keynote}[What the example is showing]
Three things, each general. First, the two factors were chosen so that neither is separately hard and the product is, which is how these problems are always built: look for the factorisation before looking for a derivative. Second, multiplication by $1/(1-x)$ is partial summation of coefficients, and multiplication by $1/(1-x)^{2}$ is partial summation twice; those two facts alone handle most rational multipliers. Third, the answer came out as $20!$ times a rational number, which is what a twentieth derivative at $0$ of any function with rational series coefficients must look like. If your answer to such a problem is not of that shape, it is wrong.
\end{keynote}

The same three moves handle the Fibonacci generating function, which is the other standard specimen. Since $\dfrac{x}{1-x-x^{2}}=\sum_{n\geq1}F_{n}x^{n}$ with $F_{n}$ the Fibonacci numbers, the twentieth derivative at $0$ is $20!\,F_{20}=20!\cdot6765$, about $1.646\times10^{22}$. The recursion $F_{n}=F_{n-1}+F_{n-2}$ is just the statement that multiplying the series by $1-x-x^{2}$ leaves $x$, read one coefficient at a time. There is no closed-form $n$-th derivative here at all, and there does not need to be.

One more, because it produces an answer worth recognising. For $g(x)=\dfrac{\eu^{x}}{1-x}$, partial summation of the exponential coefficients gives
\[
[x^{n}]g=\sum_{k=0}^{n}\frac{1}{k!},\qquad g^{(n)}(0)=n!\sum_{k=0}^{n}\frac{1}{k!} .
\]
That is an integer, and it is the nearest integer to $n!\,\eu$ for every $n\geq1$, since the tail $n!\sum_{k>n}1/k!$ is less than $1/2$. At $n=10$ the value is $9864101$, while $10!\,\eu=9864101.0991$; contour differentiation returns $9864101.0$. A problem whose answer is ``the nearest integer to $10!\,\eu$'' is one nobody solves by differentiating.

\subsection{The same derivative three ways}

Three of the four homes are open to one and the same function, and it is worth walking into all three once. Take $f(x)=x^{2}\eu^{3x}$ and $f^{(20)}(0)$.

\emph{By the closed form.} Differentiate three times and look: $f'=(2x+3x^{2})\eu^{3x}$, $f''=(2+12x+9x^{2})\eu^{3x}$, $f'''=(18+54x+27x^{2})\eu^{3x}$. The pattern is $\bigl(\alpha_{n}+\beta_{n}x+3^{n}x^{2}\bigr)\eu^{3x}$ with $\beta_{n}=2n\,3^{n-1}$ and $\alpha_{n}=n(n-1)3^{n-2}$, provable by a one-line induction. At $x=0$ only $\alpha_{n}$ survives.

\emph{By Leibniz.} Three nonzero terms, as in the block below, and only the $k=2$ term survives at $x=0$: $2\binom n2 3^{n-2}$.

\emph{By the coefficient.} $\eu^{3x}=\sum_{k\geq0}3^{k}x^{k}/k!$, so $f=x^{2}\eu^{3x}=\sum_{k\geq0}3^{k}x^{k+2}/k!$ and $[x^{20}]f=3^{18}/18!$. The bridge gives $f^{(20)}(0)=20!\cdot3^{18}/18!=20\cdot19\cdot3^{18}$.

All three give $147219785820$, and contour differentiation agrees. The comparison is instructive because the three differ in what they generalise to. The closed form is the only one that answers a question about $f^{(20)}(1)$ rather than $f^{(20)}(0)$. Leibniz is the only one that works when the second factor is unknown and only combinations of its derivatives are given. Coefficient extraction is the only one that survives when the second factor is replaced by something with no closed-form high derivative, such as $\ln(1+x)$. Choose by asking which of those three futures the problem is heading towards.

\section{Products, composites, and the sums that truncate}

Coefficient extraction assumes you can get the series. When you cannot, but the function is a product with one simple factor, the binomial theorem for derivatives does the job directly, and it does it in finitely many terms.

\begin{keynote}[Leibniz's rule and the Cauchy product are one statement]
Take the Cauchy product $[x^{n}](uv)=\sum_{j=0}^{n}[x^{j}]u\cdot[x^{n-j}]v$ and translate every coefficient into a derivative by the bridge, $[x^{m}]w=w^{(m)}(0)/m!$. Multiplying through by $n!$ turns the coefficient $n!/\bigl(j!(n-j)!\bigr)$ into $\binom nj$ and leaves
\[
(uv)^{(n)}(0)=\sum_{j=0}^{n}\binom nj u^{(j)}(0)\,v^{(n-j)}(0),
\]
which is Leibniz's rule. So the Leibniz expansion and the Maclaurin coefficient are not competitors: they are the same identity written in two systems of units, derivatives on one side and coefficients on the other. Use whichever side the data arrive in. If you are handed derivative values, use Leibniz; if you are handed series, use the product.
\end{keynote}

\tech[Leibniz rule for high derivatives of a product]{Leibniz's rule for $(fg)^{(n)}$}{medium}
\cue A product differentiated to high order, especially a polynomial times an exponential or a sinusoid; or a product where the data given are certain \emph{combinations} of one factor's derivatives rather than the factor itself.
\method Use
\[
(uv)^{(n)}=\sum_{k=0}^{n}\binom nk u^{(k)}v^{(n-k)} .
\]
When $u$ is a polynomial of degree $d$, every term with $k>d$ vanishes and the sum has $d+1$ terms regardless of $n$. When $v=\eu^{x}$, all $v^{(j)}$ are equal and the sum becomes $\eu^{x}\sum_{k}\binom nk u^{(k)}$. Choose which factor plays $u$ by asking which one runs out of derivatives first.

\begin{worked}
$f(x)=x^{2}\eu^{3x}$, and $f^{(20)}(0)$.

Take $u=x^{2}$ and $v=\eu^{3x}$. Then $u=x^{2}$, $u'=2x$, $u''=2$, and $u^{(k)}=0$ for $k\geq3$, so the sum has three terms:
\[
f^{(n)}(x)=\binom n0 x^{2}\cdot3^{n}\eu^{3x}+\binom n1 2x\cdot3^{n-1}\eu^{3x}+\binom n2 2\cdot3^{n-2}\eu^{3x}.
\]
At $x=0$ the first two terms die and only the third survives:
\[
f^{(n)}(0)=2\binom n2 3^{n-2}=n(n-1)3^{n-2},\qquad f^{(20)}(0)=20\cdot19\cdot3^{18}=147219785820 .
\]
Contour differentiation returns $1.4721979\times10^{11}$. Note that the evaluation point did half the work: the general $f^{(n)}(x)$ has three terms, and $f^{(n)}(0)$ has one.

Now the version where Leibniz is the only route, because the factor $g$ is never given. Let $f(x)=\eu^{x}g(x)$, where all that is known about $g$ is
\[
g(0)+g'''(0)=5,\qquad g'(0)+g''(0)=2 .
\]
Find $f'''(0)$. With $v=\eu^{x}$ every $v^{(j)}(0)=1$, so
\[
f'''(0)=\sum_{k=0}^{3}\binom 3k g^{(3-k)}(0)=g'''(0)+3g''(0)+3g'(0)+g(0).
\]
Regrouping into the combinations supplied, this is $\bigl(g(0)+g'''(0)\bigr)+3\bigl(g'(0)+g''(0)\bigr)=5+3\cdot2=11$. The binomial coefficients $1,3,3,1$ are symmetric, which is exactly why the two given combinations were enough; the problem was designed around that symmetry.

A third instance, at an order chosen to be absurd. Let $f(x)=x^{3}\sin x$ and find $f^{(50)}(0)$. Take $u=x^{3}$, so $u(0)=u'(0)=u''(0)=0$ and $u'''(0)=6$: of the four surviving Leibniz terms, three are killed by the evaluation point and one remains,
\[
f^{(50)}(0)=\binom{50}{3}\cdot6\cdot\sin^{(47)}(0)=19600\cdot6\cdot\sin\!\left(\frac{47\pi}{2}\right)=-117600 ,
\]
using $47\equiv3\pmod4$ so that $\sin(47\pi/2)=\sin(3\pi/2)=-1$. Two homes were used at once: Leibniz truncated the sum and the cyclic orbit evaluated the surviving factor. The fourth home confirms it independently, since $x^{3}\sin x=\sum_{m\geq0}(-1)^{m}x^{2m+4}/(2m+1)!$ gives $[x^{50}]f=-1/47!$ and hence $f^{(50)}(0)=-50!/47!=-50\cdot49\cdot48=-117600$. The same computation at $n=12$ predicts $\binom{12}{3}\cdot6\cdot\sin(9\pi/2)=1320$, and contour differentiation returns $1320.0000$.
\end{worked}

\begin{trap}
Reversing which factor carries $k$ and which carries $n-k$. With symmetric binomial coefficients this is harmless, and solvers therefore never learn to be careful, which is fatal the first time the two factors differ in an unsymmetric way. The second trap is forgetting the truncation and writing out all $n+1$ terms of a sum in which only three are nonzero.
\end{trap}

\begin{seealsob}The generalised product rule for $n$ factors; Faa di Bruno for composites; Maclaurin coefficient extraction.\end{seealsob}

Leibniz handles products. It says nothing about composites, and for good reason: there is no binomial-style formula for $\bigl(g\circ f\bigr)^{(n)}$, because the chain rule mixes orders rather than distributing them. The correct generalisation exists, is more complicated, and is mostly useful as a structural statement rather than as a computational recipe.

\tech[Faa di Bruno and Bell polynomials]{Faa di Bruno's formula for composites}{hard}
\cue A high derivative of a composite $g\bigl(f(x)\bigr)$ where Leibniz does not apply: $\eu^{\sin x}$, $\ln(1+\sin x)$, $\cos(x^{2})$, to fourth or fifth order.
\method The formula is
\[
\frac{\mathrm{d}^{n}}{\mathrm{d}x^{n}}g\bigl(f(x)\bigr)=\sum_{k=1}^{n}g^{(k)}\bigl(f(x)\bigr)\,B_{n,k}\bigl(f'(x),f''(x),\dots,f^{(n-k+1)}(x)\bigr),
\]
where $B_{n,k}$ are the partial Bell polynomials. For $n\leq4$ it is faster to differentiate directly and collect, or to compose the two Maclaurin series; the formula earns its keep for structural arguments, for deciding which terms can survive at a given point, and for organising a computation that would otherwise be an unlabelled pile of terms. Written out, the first four orders are
\begin{align*}
(g\circ f)' &= g'f', \qquad (g\circ f)'' = g''(f')^{2}+g'f'',\\
(g\circ f)''' &= g'''(f')^{3}+3g''f'f''+g'f''',\\
(g\circ f)'''' &= g''''(f')^{4}+6g'''(f')^{2}f''+g''\bigl(3(f'')^{2}+4f'f'''\bigr)+g'f'''',
\end{align*}
with every $g^{(k)}$ evaluated at $f(x)$. The coefficients $1$; $1,1$; $1,3,1$; $1,6,3{+}4,1$ are counts of set partitions, not binomial coefficients, and that is the structural difference from Leibniz: a product distributes derivatives between two slots, while a composite distributes them among an unknown number of copies of $f$.

\begin{worked}
$f(x)=\eu^{\sin x}$, and the first four derivatives at $0$.

Series composition is the practical route. With $u=\sin x=x-\tfrac{x^{3}}{6}+O(x^{5})$,
\[
\eu^{u}=1+u+\frac{u^{2}}{2}+\frac{u^{3}}{6}+\frac{u^{4}}{24}+O(u^{5}),
\]
and to degree $4$ we need $u=x-\tfrac{x^{3}}{6}$, $u^{2}=x^{2}-\tfrac{x^{4}}{3}$, $u^{3}=x^{3}$, $u^{4}=x^{4}$. Collecting,
\[
\eu^{\sin x}=1+x+\frac{x^{2}}{2}+\left(-\frac16+\frac16\right)x^{3}+\left(-\frac16+\frac1{24}\right)x^{4}+O(x^{5})
=1+x+\frac{x^{2}}{2}+0\cdot x^{3}-\frac{x^{4}}{8}+O(x^{5}).
\]
The bridge gives $f'(0)=1$, $f''(0)=1$, $f'''(0)=0$, and $f^{(4)}(0)=24\cdot\bigl(-\tfrac18\bigr)=-3$. Contour differentiation returns $-3.0000000$. The vanishing third derivative is the interesting entry, and the formula above says exactly why. At $x=0$ the inner derivatives are $f'=1$, $f''=0$, $f'''=-1$, $f''''=0$, and every $g^{(k)}$ equals $\eu^{\sin0}=1$. The third-order line therefore reads $1\cdot1+3\cdot1\cdot1\cdot0+1\cdot(-1)=0$: the $k=3$ term and the $k=1$ term cancel exactly, and the middle term was already dead because $\sin''(0)=0$. The fourth-order line reads $1+6\cdot1\cdot0+\bigl(3\cdot0+4\cdot1\cdot(-1)\bigr)+0=-3$, matching the series computation. Away from the origin nothing cancels: at $x=0.4$ the same four-term line gives $-5.6262947$, which contour differentiation confirms.

Two more of the same kind. For $\ln(1+\sin x)$, the same substitution into $\ln(1+u)=u-\tfrac{u^{2}}{2}+\tfrac{u^{3}}{3}-\tfrac{u^{4}}{4}$ gives $x-\tfrac{x^{2}}{2}+\tfrac{x^{3}}{6}-\tfrac{x^{4}}{12}$, so $f^{(4)}(0)=24\cdot\bigl(-\tfrac1{12}\bigr)=-2$, confirmed numerically as $-2.0000000$. For $\cos(x^{2})=1-\tfrac{x^{4}}{2}+O(x^{8})$, only every fourth derivative is nonzero at $0$, and $f^{(4)}(0)=24\cdot\bigl(-\tfrac12\bigr)=-12$, confirmed as $-12.0000000$.
\end{worked}

\begin{trap}
Writing $\bigl(g\circ f\bigr)^{(n)}=g^{(n)}(f)\cdot\bigl(f'\bigr)^{n}$. That is only the $k=n$ term of the formula, and every lower $k$ contributes. The error is seductive because it is correct at $n=1$ and gives a plausible-looking answer at every higher $n$. Check against $n=2$, where the truth is $g''(f)(f')^{2}+g'(f)f''$ and the missing second term is usually the larger one.
\end{trap}

\begin{seealsob}Outside-in chain-rule discipline; Maclaurin coefficient extraction; Leibniz's rule for a product.\end{seealsob}

\section{Sums of derivatives, defining equations, and non-integer orders}

Three further shapes complete the chapter. The first asks not for one derivative but for all of them added together, which sounds worse and is usually easier.

\tech[Summing all derivatives at a point]{Summing all derivatives at a point}{medium}
\cue A sum $\sum_{n\geq0}f^{(n)}(a)$, or a sum of derivatives of all orders evaluated at one point.
\method Two outcomes, and deciding which one applies is the whole technique. If $f$ is a polynomial of degree $d$, the sum truncates at $n=d$ and is a finite computation. If $f$ is self-similar under differentiation, meaning $f^{(n)}=c^{n}f$ or a scaled phase shift of $f$, the sum is geometric. For $f=\eu^{ax}$ with $\abs{a}<1$, $\sum_{n\geq0}f^{(n)}(0)=\sum_{n\geq0}a^{n}=\dfrac{1}{1-a}$. For a sinusoid, write the phase-shift form and sum a complex geometric series using $\cos\bigl(\phi+\tfrac{n\pi}{2}\bigr)=\operatorname{Re}\bigl(\eu^{\iu\phi}\iu^{\,n}\bigr)$.

\begin{worked}
$f(x)=\sqrt{1+\cos x}$, and $\displaystyle\sum_{n\geq0}f^{(n)}(\pi/2)$.

The radical is the obstruction, and the half-angle identity removes it: $1+\cos x=2\cos^{2}(x/2)$, so $f(x)=\sqrt2\,\abs{\cos(x/2)}$, and near $x=\pi/2$ the cosine is positive, so $f(x)=\sqrt2\cos(x/2)$ with no absolute value. Numerically $\sqrt{1+\cos1.3}=1.1258325$ and $\sqrt2\cos(0.65)=1.1258325$.

Now the table entry applies with $a=\tfrac12$:
\[
f^{(n)}(x)=\sqrt2\,2^{-n}\cos\!\left(\frac x2+\frac{n\pi}{2}\right),\qquad
f^{(n)}(\pi/2)=\sqrt2\,2^{-n}\cos\!\left(\frac\pi4+\frac{n\pi}{2}\right).
\]
Contour differentiation at $n=0,\dots,6$ returns $1$, $-0.5$, $-0.25$, $0.125$, $0.0625$, $-0.03125$, $-0.015625$, matching the formula exactly. Summing, with $\cos\bigl(\tfrac\pi4+\tfrac{n\pi}{2}\bigr)=\operatorname{Re}\bigl(\eu^{\iu\pi/4}\iu^{\,n}\bigr)$,
\[
\sum_{n\geq0}f^{(n)}(\pi/2)=\sqrt2\operatorname{Re}\sum_{n\geq0}\eu^{\iu\pi/4}\left(\frac{\iu}{2}\right)^{n}
=\sqrt2\operatorname{Re}\frac{\eu^{\iu\pi/4}}{1-\iu/2}
=\operatorname{Re}\frac{1+\iu}{1-\iu/2}=\operatorname{Re}\bigl(0.4+1.2\,\iu\bigr)=\frac25 .
\]
The ratio $\iu/2$ has modulus $\tfrac12<1$, so the series converges. Summing the closed-form terms numerically to $n=300$ gives $0.40000000$.

The exponential case in one line: for $f(x)=\eu^{x/3}$, $f^{(n)}(0)=3^{-n}$ and $\sum_{n\geq0}3^{-n}=\dfrac32$.

A damped oscillation, to show the complex geometric sum in its natural setting. For $f(x)=\eu^{-x/2}\cos(x/2)$ write $f=\operatorname{Re}\eu^{cx}$ with $c=\tfrac{-1+\iu}{2}$, so $f^{(n)}(0)=\operatorname{Re}\bigl(c^{n}\bigr)$ and $\abs{c}=\tfrac{1}{\sqrt2}<1$, which settles convergence before any summing. Then
\[
\sum_{n\geq0}f^{(n)}(0)=\operatorname{Re}\frac{1}{1-c}=\operatorname{Re}\frac{1}{\tfrac32-\tfrac{\iu}{2}}=\operatorname{Re}\frac{2}{3-\iu}=\operatorname{Re}\frac{3+\iu}{5}=\frac35 .
\]
Summing the terms numerically to $n=400$ gives $0.60000000$, and the individual check $f'''(0)=\operatorname{Re}(c^{3})=\tfrac14$ is confirmed by contour differentiation. Note that the damping is what makes the sum converge: for the undamped $\cos x$ the derivatives at $0$ run $1,0,-1,0,1,\dots$ and the series has no sum at all.

The polynomial case is the one people overlook, because it looks too easy to be the intended reading. For $f(x)=x^{4}-2x^{3}+7$, every derivative past the fourth is identically zero, so the sum has five terms. At $a=1$: $f(1)=6$, $f'(1)=4-6=-2$, $f''(1)=12-12=0$, $f'''(1)=24-12=12$, $f^{(4)}(1)=24$, and the total is $40$. No convergence question arises, because the series was finite from the start.
\end{worked}

\begin{trap}
Assuming that $\sum_{n\geq0}f^{(n)}(a)$ equals $f(a+1)$ by Taylor's theorem. It does not, in general: Taylor's series at $a$ with $x-a=1$ has terms $f^{(n)}(a)/n!$, not $f^{(n)}(a)$, so the two sums agree only in the special circumstances where the missing factorials cancel against the growth of the derivatives. The other check to run before summing is convergence: for $\eu^{ax}$ the sum is $1/(1-a)$ only when $\abs{a}<1$, and at $a=1$ the sum of $f^{(n)}(0)=1$ diverges outright.
\end{trap}

\begin{seealsob}$n$-th derivative closed forms; Maclaurin coefficient extraction; sum the hidden series first.\end{seealsob}

The next shape is a function that is not given by a formula at all but by a defining integral or integral representation. Differentiating such a thing repeatedly by brute force is out of the question. Converting it to a differential equation is not.

\tech{Special functions via the defining ODE recursion}{hard}
\cue A function defined by an integral or an integral representation, with a specific $f^{(n)}(0)$ wanted: Dawson's function $F(x)=\eu^{-x^{2}}\int_{0}^{x}\eu^{t^{2}}\dd t$, or Bessel's $J_{0}(x)=\frac1\pi\int_{0}^{\pi}\cos(x\sin\theta)\dd\theta$.
\method Differentiate the definition once, using the product rule and the fundamental theorem of calculus, to obtain a differential equation with no integral in it. Then feed a power series into the equation and match coefficients, which produces a recursion; parity halves the work. When the representation is an integral over a parameter, differentiating under the integral sign and evaluating a Wallis integral is often faster still.

\begin{worked}
Dawson's function $F(x)=\eu^{-x^{2}}\displaystyle\int_{0}^{x}\eu^{t^{2}}\dd t$, and $F^{(5)}(0)$.

Differentiate the product. The exponential gives $-2x\eu^{-x^{2}}\int_{0}^{x}\eu^{t^{2}}\dd t=-2xF(x)$, and the fundamental theorem gives $\eu^{-x^{2}}\cdot\eu^{x^{2}}=1$. So
\[
F'=1-2xF,\qquad F(0)=0,
\]
and the integral has disappeared. Now substitute $F=\sum_{k\geq0}a_{k}x^{k}$. Matching coefficients of $x^{m}$ gives $(m+1)a_{m+1}=[m=0]-2a_{m-1}$, that is $a_{1}=1$ and $a_{m+1}=-\dfrac{2a_{m-1}}{m+1}$ for $m\geq1$. Since $F$ is odd (the integrand is even, so the integral is odd, and $\eu^{-x^{2}}$ is even), all even coefficients vanish and only the odd chain matters:
\[
a_{1}=1,\quad a_{3}=-\frac{2}{3},\quad a_{5}=-\frac{2a_{3}}{5}=\frac{4}{15},\quad a_{7}=-\frac{2a_{5}}{7}=-\frac{8}{105}.
\]
Hence $F(x)=x-\tfrac23x^{3}+\tfrac4{15}x^{5}-\tfrac8{105}x^{7}+\cdots$, and the bridge gives
\[
F^{(5)}(0)=5!\cdot\frac{4}{15}=32 .
\]
Numerical quadrature of the defining integral at $x=0.3$ gives $F(0.3)=0.2826317$ against the truncated series value $0.2826313$, the discrepancy being the $x^{9}$ term.

Bessel's $J_{0}$ takes the other route. Differentiating under the integral sign four times brings down $\sin^{4}\theta$ with sign $(+1)$:
\[
J_{0}^{(4)}(0)=\frac1\pi\int_{0}^{\pi}\sin^{4}\theta\dd\theta=\frac1\pi\cdot\frac{3\pi}{8}=\frac38 ,
\]
using the Wallis value $\int_{0}^{\pi}\sin^{4}=\tfrac{3\pi}{8}$. The check is independent: the series $J_{0}(x)=\sum_{k\geq0}(-1)^{k}(x/2)^{2k}/(k!)^{2}$ has $[x^{4}]=\tfrac{1}{64}$, so the bridge gives $4!/64=\tfrac38$ as well, and contour differentiation of the integral representation on $\abs{z}=1$ returns $0.3750000$.
\end{worked}

\begin{trap}
Trying to differentiate the integral representation five times by brute force, which reintroduces the integral at every step. Differentiate once, convert to an equation, and never look at the integral again. The second waste is computing the even coefficients of an odd function: the parity observation should come before the recursion, not after.
\end{trap}

\begin{seealsob}Differentiating under the integral sign; Maclaurin coefficient extraction; the logarithmic derivative as an operator.\end{seealsob}

The last shape is a curiosity that is nonetheless completely standard once the right basis is chosen. If differentiation is multiplication by $c$ on the exponential $\eu^{cx}$, there is no obstacle to multiplying by $c^{1/2}$ instead, and the whole theory of fractional derivatives follows from taking that remark seriously.

\tech[Fractional derivative at zero]{Fractional derivatives via the exponential basis}{hard}
\cue A differentiation order that is not an integer: $\frac{\mathrm{d}^{1/2}}{\mathrm{d}x^{1/2}}$, or an order named as a parameter $\alpha$.
\method Work on the basis where differentiation is diagonal. On exponentials, $\frac{\mathrm{d}^{\alpha}}{\mathrm{d}x^{\alpha}}\eu^{cx}=c^{\alpha}\eu^{cx}$, with the Liouville convention whose lower limit is $-\infty$; decompose the function into exponentials using $\sin x=\dfrac{\eu^{\iu x}-\eu^{-\iu x}}{2\iu}$ and $a^{x}=\eu^{x\ln a}$. On monomials the corresponding rule is
\[
\frac{\mathrm{d}^{\alpha}}{\mathrm{d}x^{\alpha}}x^{p}=\frac{\Gamma(p+1)}{\Gamma(p-\alpha+1)}\,x^{p-\alpha},
\]
which reproduces the integer rule because $\Gamma(m+1)=m!$.

\begin{worked}
$f(x)=\pi^{x}\sin x$, and $\frac{\mathrm{d}^{1/2}}{\mathrm{d}x^{1/2}}f$ evaluated at $x=0$.

Write both factors as exponentials. With $L=\ln\pi$,
\[
f(x)=\eu^{Lx}\cdot\frac{\eu^{\iu x}-\eu^{-\iu x}}{2\iu}=\frac{\eu^{(L+\iu)x}-\eu^{(L-\iu)x}}{2\iu},
\]
a difference of two pure exponentials, each with positive real part in the exponent, so the Liouville integral converges. Half-differentiating multiplies each by the square root of its own constant, and at $x=0$ both exponentials equal $1$:
\[
\left.\frac{\mathrm{d}^{1/2}}{\mathrm{d}x^{1/2}}f\right|_{0}=\frac{(L+\iu)^{1/2}-(L-\iu)^{1/2}}{2\iu}=\operatorname{Im}\bigl[(L+\iu)^{1/2}\bigr],
\]
the last step because the two constants are complex conjugates and $\bigl(\bar z\bigr)^{1/2}=\overline{z^{1/2}}$ on the principal branch. With $L=1.1447299$, the modulus of $L+\iu$ is $1.520002$ and its argument is $0.718019$, so the principal square root has modulus $1.232884$ and argument $0.359009$, giving
\[
\left.\frac{\mathrm{d}^{1/2}}{\mathrm{d}x^{1/2}}f\right|_{0}=1.232884\sin(0.359009)=0.4331699 .
\]
Evaluating the Liouville integral $\frac{1}{\Gamma(1/2)}\dvop{x}\int_{0}^{\infty}s^{-1/2}f(x-s)\dd s$ numerically at $x=0$ returns $0.4331699$.

The monomial rule deserves one check of its own, because it is the fastest way to see that the definition is the right one. Half-differentiating $x$ gives $\dfrac{\Gamma(2)}{\Gamma(3/2)}x^{1/2}=\dfrac{2}{\sqrt\pi}x^{1/2}$, and half-differentiating that again gives $\dfrac{2}{\sqrt\pi}\cdot\dfrac{\Gamma(3/2)}{\Gamma(1)}x^{0}=\dfrac{2}{\sqrt\pi}\cdot\dfrac{\sqrt\pi}{2}=1$, which is $\dvop{x}x$. Two half-derivatives compose to one whole derivative, exactly as the notation promises. Similarly $\frac{\mathrm{d}^{1/2}}{\mathrm{d}x^{1/2}}x^{3/2}=\dfrac{\Gamma(5/2)}{\Gamma(2)}x=\dfrac{3\sqrt\pi}{4}x=1.3293404\,x$.
\end{worked}

\begin{trap}
Choosing the wrong branch of the complex square root, which flips the sign of the answer. Stay on the principal branch, where the argument lies in $(-\pi,\pi]$, and check that conjugate inputs produce conjugate outputs. The other trap is mixing conventions: the Liouville form used here has lower limit $-\infty$ and makes exponentials eigenfunctions, while the Riemann-Liouville form with lower limit $0$ and the Caputo form do not, and the three disagree by terms that depend on the initial values. State which one you are using.
\end{trap}

\begin{seealsob}Complex exponentials for high derivatives; $n$-th derivative closed forms; Maclaurin coefficient extraction.\end{seealsob}

\subsection{A fiftieth derivative, in full}

One closing problem, chosen because it needs three of the four homes in sequence and is hopeless without them. Find $f^{(50)}(0)$ for
\[
f(x)=\frac{\sin\bigl(x^{2}\bigr)}{1-x^{2}} .
\]

\emph{First move: the variable appears only as $x^{2}$.} Substitute $u=x^{2}$, so $f(x)=g(u)$ with $g(u)=\dfrac{\sin u}{1-u}$. Every odd derivative at the origin vanishes, and $f^{(50)}(0)=50!\,[u^{25}]g$. A fiftieth derivative has become a twenty-fifth coefficient, and half the possible answers have been deleted without a computation. The check runs the other way too: $f^{(49)}(0)=0$, and contour differentiation returns $-5\times10^{52}$ against a scale of $10^{64}$, which is zero to the precision of the quadrature.

\emph{Second move: the multiplier $1/(1-u)$ is partial summation.} The coefficients of $g$ are the partial sums of the coefficients of $\sin u$, so
\[
[u^{25}]g=\sum_{\substack{j\leq25\\ j\text{ odd}}}[u^{j}]\sin u=\sum_{m=0}^{12}\frac{(-1)^{m}}{(2m+1)!}
=1-\frac{1}{3!}+\frac{1}{5!}-\cdots+\frac{1}{25!} .
\]

\emph{Third move: recognise the partial sum.} That is the Maclaurin series for $\sin1$, truncated after thirteen terms. The first omitted term is $1/27!\approx9.2\times10^{-29}$, so the partial sum agrees with $\sin1=0.8414709848078965$ in every digit a calculator will print. Exactly,
\[
[u^{25}]g=\frac{13052233190723859821607001}{25!},
\]
and therefore
\[
f^{(50)}(0)=\frac{50!}{25!}\cdot13052233190723859821607001=2.5592577\times10^{64},
\]
an integer because $25!$ divides $50!$. Contour differentiation on a circle of radius $0.75$ returns $2.5592577\times10^{64}$, and the magnitude estimate agrees as well: the nearest singularities of $f$ are the poles at $x=\pm1$, so $R=1$ and the answer should be a modest fraction of $50!=3.04\times10^{64}$, which it is.

What was never done is any differentiation at all. The three moves were a substitution, a partial summation, and a recognition, and each of them is one line. That is what the whole chapter has been arguing: a request for a fiftieth derivative is a request to locate the pattern, and the pattern is always in one of four places.

\section{Recognition matrix}
\begin{recog}{@{}p{0.31\textwidth}p{0.35\textwidth}p{0.28\textwidth}@{}}
\rowcolor{mist}\mh{If you see} & \mh{Reach for} & \mh{If that fails} \\
\midrule
\endhead
$f^{(n)}$ wanted symbolically in $n$ & Three orders, conjecture, induction & Reduce $f$ by an identity first \\
$f^{(20)}(0)$, $f^{(100)}(0)$, $f^{(2026)}(0)$ & $f^{(n)}(0)=n!\,[x^{n}]f$ & Parity first: the answer may be $0$ \\
An answer divided by $n!$ or $\Gamma(n)$ & Coefficient extraction; the factorials cancel & Read the normalisation as the hint \\
$\sin^{2}x$, $\cos^{3}x$, $\sin x\cos x$ & Linearise, then $a^{n}\sin(ax+b+n\pi/2)$ & Complex exponentials \\
A rational function, high order & Partial fractions, then $(-1)^{n}n!a^{n}(ax+b)^{-n-1}$ & Geometric series coefficients \\
$\ln(ax+b)$ to order $n$ & $(-1)^{n-1}(n-1)!\,a^{n}(ax+b)^{-n}$ & Differentiate once, then the previous row \\
$\eu^{ax}\sin bx$ or $\eu^{ax}\cos bx$ & $\operatorname{Im}$ or $\operatorname{Re}$ of $(a+\iu b)^{n}$ & Repeated squaring for a specific $n$ \\
Odd times even, evaluated at $0$ & Parity: half the derivatives vanish & Shift to the centre of symmetry first \\
Polynomial times anything & Leibniz; the sum stops at $\deg$ & Series product instead \\
Only combinations of $g^{(k)}(0)$ given & Leibniz, then regroup into those combinations & Check the binomial symmetry \\
$g(f(x))$ to order $4$ or $5$ & Compose the two Maclaurin series & Faa di Bruno with Bell polynomials \\
$\sum_{n\geq0}f^{(n)}(a)$ & Polynomial: truncates. Self-similar: geometric & Phase-shift form plus $\operatorname{Re}$ \\
$f$ defined by an integral & Differentiate once into an ODE, then recurse & Differentiate under the sign \\
An order that is not an integer & $\frac{\mathrm{d}^{\alpha}}{\mathrm{d}x^{\alpha}}\eu^{cx}=c^{\alpha}\eu^{cx}$ & $\Gamma(p+1)/\Gamma(p-\alpha+1)$ on monomials \\
$\ln(1+x)$ or $\arctan x$ as a factor & Cauchy product with the other factor's series & Times $1/(1-x)$ is partial summation \\
A generating function with a recursion & Read $[x^{n}]$ off the recursion, then times $n!$ & Partial fractions if the denominator factors \\
An answer that is not $n!$ times a rational & Suspect an arithmetic slip & Recheck the factorial at the last step \\
$f(x)=g(x^{k})$, derivative at $0$ & Zero unless $k\mid n$; else $n!\,[x^{n/k}]g$ & Substitute $t=x^{k}$ first \\
A repeated factor in the denominator & $(-1)^{n}\dfrac{(m+n-1)!}{(m-1)!}(x-c)^{-m-n}$ & The worst pole sets the growth \\
$\eu^{ax}\sin bx$ at general $x$ & $r^{n}\eu^{ax}\sin(bx+n\theta)$, $r=\sqrt{a^{2}+b^{2}}$ & Set $x=0$ and take $\operatorname{Im}(a+\iu b)^{n}$ \\
An answer of the wrong order of magnitude & Compare with $n!/R^{n}$, $R$ to the nearest pole & Recount the poles of $f$ \\
$(1-4x)^{-1/2}$ or a surd factor & Binomial series: $[x^{n}]=\binom{2n}{n}$ & $\binom{\alpha}{n}$ in general \\
\end{recog}
